\documentclass[11pt,a4paper]{article}
\usepackage[T1]{fontenc}
\usepackage{isabelle,isabellesym}
\usepackage{amsmath,amssymb}

% Allow discretionary line breaks inside long typewriter identifiers
% (e.g. MSOinHOL_lowenheim_skolem_lemmas) without changing their rendering.
% Encourage breaking (penalty < hyphenpenalty so TeX actually uses it).
\newcommand{\bk}{\_\penalty10\hspace{0pt}}

% Extra stretch lets TeX absorb the leftover whitespace once it does break,
% so it stops preferring an overfull line.
\setlength{\emergencystretch}{3em}

% this should be the last package used
\usepackage{pdfsetup}

% urls in roman style, theory text in italics
\urlstyle{rm}
\isabellestyle{it}

\begin{document}

\title{Monadic Second-Order Logic in HOL:\\
       Deep and Shallow Embeddings with Automated Faithfulness\\
       (Isabelle/HOL dataset)}
\author{Christoph Benzm\"uller \and Daniel Kirchner}
\maketitle

\begin{abstract}
We develop, in Isabelle/HOL, three embeddings of monadic second-order logic
(MSO) into classical higher-order logic side by side: a \emph{deep} embedding
(an inductive datatype with an explicitly defined satisfaction relation), a
\emph{maximal-shallow} embedding, and a \emph{minimal-shallow} embedding --- the
last a locale parametrised by an interpretation and by first- and second-order
assignments. The enabling ingredient is a \emph{two-sorted} capture-avoiding
substitution apparatus (predication, substitution, alphabetic renaming, and a
substitution lemma, developed once for each of the two binder namespaces), in
which each binder is transparent for the other. Faithfulness of all three
embeddings is mechanised and automated.

The central result is a fully mechanised two-sorted downward
L\"owenheim--Skolem theorem. Its corollary identifies range-relative validity
with the general (Henkin-style) reading of MSO, while comprehension witnesses
that the standard reading is strictly stronger; an elementary-substructure
refinement recovers the standard reading from the minimal embedding as well, so
the two readings become one device under two interpretation classes (all
interpretations versus the elementary ones). We further exercise the embeddings
on classical MSO landmarks: Boolean closure and graph operations hold under the
standard (full second-order) reading and fail under the general one, making the
dichotomy concrete, while reachability and $2$-colorability are non-theorems,
refuted throughout with \texttt{nitpick} countermodels.
\end{abstract}

% List sections and subsections only.  The experiment theories carry many
% subsubsection headings; showing them would make the Experiments entry
% dominate the table of contents.
\setcounter{tocdepth}{2}
\tableofcontents
\newpage

\section{Overview}

This entry extends the deep-and-shallow embedding methodology of
\emph{Faithful Logic Embeddings in HOL --- Deep and Shallow}~\cite{C98} from
propositional modal logic to a logic with genuine quantifier binding. The shallow-embedding approach goes
back to the modal-logic embeddings of Benzm\"uller and Paulson~\cite{J21,J23}
and underlies the LogiKEy methodology~\cite{J48}. The background on MSO and its
model theory follows the standard references~\cite{Thomas1997,Courcelle2012,Buchi1960};
the notions of general (Henkin) and standard second-order semantics follow
\cite{henkin1950,Shapiro1991,Vaananen2011,Manzano1996}. Locales~\cite{ballarin2014}
parametrise the minimal embedding throughout. Within the AFP, MSO on words has
been treated algorithmically by Traytel and Nipkow~\cite{MSO_Regex_Equivalence-AFP},
who verify decision procedures based on derivatives of regular expressions; the
present entry is complementary, being concerned not with decision procedures but
with faithful embeddings of MSO into HOL and with the resulting model theory
(the general versus standard readings and a two-sorted L\"owenheim--Skolem
theorem).

The development is organised as follows.

\paragraph{Syntax and the deep embedding.}
\texttt{MSOinHOL\_preliminaries} fixes the variable namespaces, domains, and
assignment-update operations. \texttt{MSOinHOL\_deep} introduces MSO formulae as
an inductive datatype together with the deep satisfaction relation and the
global validity notions $\models^{d}$ (general; all interpretations) and
$\models^{d'}$ (standard; full second-order domain).

\paragraph{Shallow embeddings.}
\texttt{MSOinHOL\_shallow} gives the maximal-shallow embedding.
\texttt{MSOinHOL\bk shallow\bk minimal\bk locale} and
\texttt{MSOinHOL\bk shallow\bk minimal} give the minimal-shallow embedding as a
locale \texttt{MinS} parametrised by an interpretation and two assignments,
together with its translation into HOL.

\paragraph{Substitution.}
\texttt{MSOinHOL\bk deep\bk subst\bk lemma} develops the two-sorted
capture-avoiding substitution and renaming machinery and the substitution
lemmas \texttt{SubstitutionLemma} and \texttt{SubstitutionLemma2}, one per
binder namespace; \texttt{MSOinHOL\bk subst\bk extras} collects auxiliary
renaming-preserves-truth results.

\paragraph{Faithfulness.}
\texttt{MSOinHOL\bk faithfulness\bk locale} proves the faithfulness theorems
relating the deep, maximal-shallow and minimal-shallow embeddings
(\texttt{FaithfulSD}, \texttt{FaithfulMD}, \texttt{FaithfulMS}, and
\texttt{FaithfulMS\_all}, which identifies minimal validity with range-relative
deep validity).

\paragraph{Comprehension and L\"owenheim--Skolem.}
\texttt{MSOinHOL\bk comprehension} isolates the comprehension schema that
separates the readings. \texttt{MSOinHOL\bk lowenheim\bk skolem\bk lemmas}
builds the two-sorted Tarski--Vaught stage construction, and
\texttt{MSOinHOL\bk lowenheim\bk skolem} assembles the headline results: the
downward L\"owenheim--Skolem theorem \texttt{DownwardLowenheimSkolem}, the
elementary-substructure relation \texttt{ElementarySubstructure}, the
faithfulness corollaries \texttt{Faithful\bk to\bk Henkin} and
\texttt{Faithful\bk to\bk Standard}, and the limitative result
\texttt{Standard\bk strictly\bk stronger}.

\paragraph{Experiments.}
The \texttt{experiments} theories exercise classical MSO landmarks across the
embeddings --- Boolean closure and graph operations hold under the standard
reading and fail under the general one, while reachability~\cite{BasinKlarlund1995}
and $2$-colorability~\cite{Thomas1997} are non-theorems, refuted throughout ---
using \texttt{nitpick} for the
distinguishing countermodels; the \texttt{\_elementary} theories carry the same
landmarks over the elementary minimal interpretation.

\paragraph{Disclosure on the use of generative AI.}
The authors used generative-AI assistants (Anthropic's Claude family of
models) as a drafting aid: to help condense prose, to suggest and shorten
selected proofs (in particular parts of the two-sorted
L\"owenheim--Skolem construction), and to help keep the \LaTeX{} sources of
this entry tidy. All mathematical content, the theory development, and the
design decisions are the authors' own. Every definition, proof, and piece
of text in the final entry has been checked by the authors, who take full
responsibility for the content.

% Generated theory content. Section/subsection headings issued here (before
% each \input) so they precede each theory's "theory ... imports ... begin"
% header in the rendering.
\section{Preliminaries}
\input{MSOinHOL_preliminaries.tex}

\section{The deep embedding}
\input{MSOinHOL_deep.tex}

\section{Two-sorted substitution}
\input{MSOinHOL_deep_subst_lemma.tex}
\subsection{Auxiliary renaming results}
\input{MSOinHOL_subst_extras.tex}

\section{Shallow embeddings}
\subsection{Maximal-shallow embedding}
\input{MSOinHOL_shallow.tex}
\subsection{Minimal-shallow embedding (locale)}
\input{MSOinHOL_shallow_minimal_locale.tex}
\subsection{Minimal-shallow embedding (translation)}
\input{MSOinHOL_shallow_minimal.tex}

\section{Faithfulness}
\subsection{The faithfulness locale}
\input{MSOinHOL_faithfulness_locale.tex}
\subsection{Faithfulness theorems}
\input{MSOinHOL_faithfulness.tex}

\section{Comprehension}
\input{MSOinHOL_comprehension.tex}

\section{Downward L\"owenheim--Skolem}
\subsection{Tarski--Vaught stage lemmas}
\input{MSOinHOL_lowenheim_skolem_lemmas.tex}
\subsection{The theorem and its corollaries}
\input{MSOinHOL_lowenheim_skolem.tex}

\section{Experiments}
\subsection{Basic landmarks}
\input{MSOinHOL_experiments.tex}
\subsection{Classical MSO landmarks}
\input{MSOinHOL_experiments_classic.tex}
\subsection{Locale-based variants}
\input{MSOinHOL_experiments_locale.tex}
\subsection{Elementary minimal embedding (setup)}
\input{MSOinHOL_shallow_minimal_elementary.tex}
\subsection{Elementary minimal embedding (landmarks)}
\input{MSOinHOL_experiments_classic_elementary.tex}
\subsection{Additional landmarks}
\input{MSOinHOL_experiments_extra.tex}

\bibliographystyle{abbrv}
\bibliography{root}

\end{document}
